% !PRINTED 146
\weberpage{146}
%<<CONT
selben Nebengruppe, etwa in $Q\pi_1$, vorkommen, und folglich giebt es zwei
Permutationen $\varkappa_1, \varkappa_2$ in $Q$, die der Bedingung

\[ \varkappa_1 \pi_1 = (0, 1), \quad \varkappa_2 \pi_1 = (0, 2) \]

genügen. Es ist daher

\[ \varkappa_1 \pi_1 \pi_1^{-1} \varkappa_2^{-1} = \varkappa_1 \varkappa_2^{-1} = (0,1)(0,2) = (0,1,2) \]

in $Q$ enthalten. Dies aber widerspricht unserer Voraussetzung, dass $Q$ keinen
dreigliedrigen Cyklus enthalten soll. Demnach kann $j$ nicht $< n$ sein. Ist aber
$j = n$, also der Grad von $Q$ gleich $\Pi(n-1)$, so müssen die $n-1$ Permutationen
(5) in $n-1$ verschiedenen Nebengruppen vorkommen, und $P$ lässt sich so darstellen:

\weq{6}{P = Q + Q(0,1) + Q(0,2) + \cdots + Q(0, n-1).}

Betrachten wir irgend eine andere Transposition, z. B. $(2,3)$, so kann diese
nicht in $Q$ und nicht in $Q(0,2)$ oder in $Q(0,3)$ vorkommen, weil sonst
$(2,3)(0,2) = (0,2,3)$ oder $(2,3)(0,3) = (0,3,2)$ in $Q$ vorkäme. Wir können also,
ohne die Allgemeinheit zu beeinträchtigen, annehmen, dass $(2,3)$ in $Q(0,1)$
vorkommt, und daraus ergiebt sich:

\gesperrt{$\beta$) Die Gruppe $Q$ enthält das Transpositionspaar}

\weq{7}{(0,1)(2,3).}

Die Gruppe $Q$ hat einen Theiler $Q_0$, der aus allen den Permutationen besteht,
die die Ziffer 0 an ihrer Stelle lassen. Da $Q$ transitiv ist, so ist, wenn
$\varkappa_1, \varkappa_2, \ldots, \varkappa_{n-1}$ Elemente aus $Q$ sind, die die 0
in 1, in 2, \ldots, in $n-1$ überführen,

\weq{8}{Q = Q_0 + Q_0 \varkappa_1 + \cdots + Q_0 \varkappa_{n-1},}

und da $Q$ vom Grade $\Pi(n-1)$ ist, so folgt hieraus, dass $Q_0$ vom Grade

\weq{9}{g = \frac{\Pi(n-1)}{n}}

ist. Da $g$ eine ganze Zahl, also $\Pi(n-1)$ durch $n$ theilbar sein muss, so
schliessen wir zunächst, dass der Fall $j = n$ niemals eintreten kann, wenn $n$
eine Primzahl ist, und für diesen Fall ist also unser Theorem II. bewiesen.

Im Allgemeinen können wir aber schliessen:

\gesperrt{$\gamma$) Ist $n$ nicht $= 4$, so sind die $n-1$ Ziffern
$1, 2, \ldots, n-1$ durch $Q_0$ noch transitiv verbunden.}
